% Round 20. PI20. Author: supported20. Independent review required.
\section{A finite correspondence error and a surface certificate}
\label{r20pi:section}

\paragraph{Scope.}
This section replaces a universal component comparison hypothesis by a
single, chosen, numerically trivial matrix correspondence.  It gives an
exact presentation of AK19's defect in terms of that correspondence.
For a native semistable surface model the error is a curve on a product
of two surfaces, so the frozen SP19 curve-comparison theorem makes its
geometric cohomology class zero.  An inverse can then be certified by a
finite rational linear system in actual divisor representatives.
The general-dimensional error is not proved homologically zero.

Throughout, $R$ is a complete discrete valuation ring with finite residue
field $k$, $\mathcal Y/R$ is an integral regular flat projective reduced
strictly semistable model of relative dimension $d\ge1$, and its generic
fiber is smooth projective and geometrically integral.  Write
$D=\coprod_iD_i$ for the disjoint union of its native smooth proper
components.  All intersections and connected pieces are retained,
including genuine higher strata.  We initially assume the $D_i$ are
geometrically integral; the final paragraph explains finite constant
extensions.  Let $\ell\ne\operatorname{char}k$ and
$\Lambda=\mathbf Q_\ell$.

\subsection{The incidence map as an actual correspondence}

Use the frozen signs
\[
 (\rho v)_{i\to j}=v_j|_{E}-v_i|_{E},\qquad
 \sigma w=(+j_{i*}w,-j_{j*}w),\qquad A=\sigma\rho.
\]
For a degree-$s$ correspondence from $D_j$ to $D_i$ we use
$\operatorname{CH}^{d+s}(D_j\times_kD_i)_{\mathbf Q}$.
It acts by pullback, intersection and proper pushforward on Chow
groups and on geometric cohomology, increasing codimension by $s$.

For $i\ne j$, let $E_{ij}$ denote the full native intersection, with its
actual closed immersions into $D_i$ and $D_j$, and put
\[
 \begin{split}
 A_{ij}&=(j_{j},j_i)_*[E_{ij}]
        \in\operatorname{CH}^{d+1}(D_j\times D_i)_{\mathbf Q},\\
 A_{ii}&=(\Delta_i)_*c_1(N_{D_i/\mathcal Y})
       =-(\Delta_i)_*\sum_{j\ne i}[E_{ij}].
 \end{split}
 \tag{PI20.1}\label{r20pi:actualA}
\]
An empty intersection contributes zero and a disconnected one
contributes the sum of its native pieces.  The last equality holds
in the actual divisor Chow group of $D_i$: the reduced total fiber
is principal, so its restriction has trivial line bundle.

The action of the matrix $A=(A_{ij})$ on actual component Chow
classes is precisely $\sigma\rho$.  Indeed an off-diagonal entry is
the indicated restriction followed by Gysin, whereas the diagonal
entry is minus the sum of the boundary-divisor self contributions.
Thus it gives simultaneously the full cohomological maps $A^H_u$
and the numerical maps $A^N_u$ of AK19/CR19.  No passage from a
numerical class to a cohomology class occurs in this construction.
The relation (PI20.1) includes every boundary component; higher
intersections do not require deleting any summand.

For later reference, the composition of actual correspondences is
\[
 \beta\circ\alpha
   =(p_{13})_*\bigl(p_{12}^*\alpha\cdot p_{23}^*\beta\bigr)
 \tag{PI20.2}\label{r20pi:composition}
\]
on the appropriate product of smooth proper strata.  Numerical
equivalence is adequate for these operations.  Every intersection
number below is a native degree on a smooth proper $k$-variety,
never a degree on the absolute $\operatorname{CH}_0(\mathcal Y)$.

\subsection{A numerical generalized inverse}

\begin{lemma}[Existence at the correspondence level]
\label{r20pi:numericalinverse}
There is a matrix of actual rational correspondences
\[
 G_{ij}\in\operatorname{CH}^{d-1}(D_j\times D_i)_{\mathbf Q}
 \tag{PI20.3}\label{r20pi:G}
\]
such that
\[
                         AGA=A
       \quad\hbox{modulo numerical equivalence of correspondences}.
 \tag{PI20.4}\label{r20pi:numidentity}
\]
In particular
\[
                         E=A-AGA
 \tag{PI20.5}\label{r20pi:error}
\]
is a numerically trivial actual matrix correspondence, each of whose
entries has codimension $d+1$ on $D_j\times D_i$.
\end{lemma}

\begin{proof}
Take rational coefficients in the category of numerical motives and
put $M=\bigoplus_i h(D_i)$.  With Jannsen's explicit convention
\[
 \operatorname{Hom}((X,p,m),(Y,q,n))
       =q\,N^{\dim X-m+n}(X\times Y)_{\mathbf Q}\,p,
\]
the incidence correspondence is a morphism $A:M\to M(1)$.
Jannsen's Theorem~1 makes this category semisimple abelian without
assuming a standard conjecture.  Choose splittings
$M=\ker A\oplus M'$ and $M(1)=\operatorname{im}A\oplus N'$.
The restriction $A:M'\to\operatorname{im}A$ is an isomorphism.
Define the numerical $G:M(1)\to M$ to be its inverse on
$\operatorname{im}A$ and zero on $N'$.  This gives (PI20.4).

The indicated morphism group consists of matrices in
$N^{d-1}(D_j\times D_i)_{\mathbf Q}$, each of which has an
actual rational Chow representative.  Choose such representatives
to obtain (PI20.3).  Their actual compositions need not satisfy
(PI20.4) in Chow or in cohomology; only the displayed numerical
identity has been proved at this point.  The composition degrees
in $AGA$ add to $1-1+1=1$, giving (PI20.5) in the stated Chow
groups.
\end{proof}

\paragraph{The precise comparison input.}
For a chosen $G$, the sufficient cohomological input is
\[
 \operatorname{cl}_{\ell}(E_{ij})=0
 \quad\hbox{in }
 H^{2d+2}((D_j\times D_i)_{\bar k},\Lambda(d+1))
 \quad\hbox{for every }i,j.
 \tag{PI20.6}\label{r20pi:errorzero}
\]
This is a finite comparison statement about the actual cycles
(PI20.5).  It is stronger than needed for the actual generic-cycle
image.  It does not follow from the action of $E$ on numerical
groups, and it is not asserted for arbitrary-dimensional models.

\subsection{An exact presentation of the actual lifting defect}

Retain AK19's canonical actual-span projection
\[
 \begin{gathered}
 \mathscr A_C^u=
 \operatorname{span}_{\Lambda}\operatorname{cl}_{\ell}
      \left(\bigoplus_i\operatorname{CH}^u(D_i)_{\mathbf Q}\right),
 \qquad p_u:\mathscr A_C^u\twoheadrightarrow C^u\otimes\Lambda,\\
 \mathscr T_C^u=\ker p_u,\qquad
 \Delta_{A,u,\ell}
 =\frac{\mathscr T_C^{u+1}\cap A^H_u\mathscr A_C^u}
              {A^H_u\mathscr T_C^u}.
 \end{gathered}
 \tag{PI20.7}\label{r20pi:delta}
\]
Actual correspondences preserve these actual spans and commute with
their numerical projections.  In particular $G^H$ preserves the
corresponding kernels $\mathscr T$ with the degree decreased by one.

\begin{theorem}[The single-error presentation]
\label{r20pi:exacterror}
For every $u$ in the component degree range, the inclusion in the
numerator of (PI20.7) induces a canonical isomorphism, relative to
the chosen actual $G$,
\[
 \boxed{\quad
 \Delta_{A,u,\ell}\simeq
 \frac{E^H_u\mathscr A_C^u}
 {E^H_u\mathscr A_C^u\cap A^H_u\mathscr T_C^u}.
 \quad}
 \tag{PI20.8}\label{r20pi:errorpresentation}
\]
Consequently
\[
 \dim_\Lambda\Delta_{A,u,\ell}
 \le\operatorname{rank}\bigl(E^H_u|_{\mathscr A_C^u}\bigr)
 \le\operatorname{rank}\bigl(E^H_u|_{H_C^u}\bigr).
 \tag{PI20.9}\label{r20pi:errorrank}
\]
If (PI20.6) holds, then $\Delta_{A,u,\ell}=0$ simultaneously
in every degree, and hence all the corresponding CR19 supported
defects vanish.
\end{theorem}

\begin{proof}
Because $E$ is numerically trivial as a correspondence, its action
on every numerical group is zero.  Hence
$E^H_u\mathscr A_C^u\subset\mathscr T_C^{u+1}$.
The identity $E=A(1-GA)$ gives the further inclusion
\[
 E^H_u\mathscr A_C^u
       \subseteq A^H_u\mathscr A_C^u,
\]
since $1-GA$ preserves actual algebraic classes.  The inclusion
therefore defines a map to $\Delta_{A,u,\ell}$ with the kernel
displayed in (PI20.8).

For surjectivity take $\alpha=A^H_ux\in\mathscr T_C^{u+1}$,
where $x\in\mathscr A_C^u$.  The class $G^H\alpha$ lies in
$\mathscr T_C^u$, since $G$ is an actual correspondence and
preserves numerical zero.  Thus
\[
 \alpha-E^H_ux=A^H_uG^H\alpha\in A^H_u\mathscr T_C^u.
 \tag{PI20.10}\label{r20pi:erroronto}
\]
So every quotient class is represented by $E^H_ux$.
This proves (PI20.8), and taking dimensions gives (PI20.9).
Under (PI20.6) the action of $E$ on all geometric cohomology is
zero, so the presentation gives every stated vanishing.
The map onto CR19's dual defect is the frozen exact sequence
(CR19.9).
\end{proof}

\begin{corollary}[An actual corrected kernel lift]
\label{r20pi:correctedlift}
Assume (PI20.6).  If $v\in\ker A^N_u$ and
$x\in\mathscr A_C^u$ is any actual-span lift, then
\[
                         x^{\flat}=x-G^HA^H_ux
 \tag{PI20.11}\label{r20pi:lift}
\]
is an actual-span lift of the same $v$ and belongs to $\ker A^H_u$.
Also, every actual algebraic class in the full cohomological
incidence image has an actual algebraic preimage:
\[
 \mathscr A_C^{u+1}\cap\operatorname{im}A^H_u
                        =A^H_u\mathscr A_C^u.
 \tag{PI20.12}\label{r20pi:imageexact}
\]
\end{corollary}

\begin{proof}
The class $A^H_ux$ is numerically zero because $A^N_uv=0$.
Therefore $G^HA^H_ux$ is numerically zero, proving
$p_u(x^{\flat})=v$.  Meanwhile
$A^H_ux^{\flat}=(A^H- A^HG^HA^H)x=E^Hx=0$.
For the second assertion write an actual class
$\alpha=A^H_u\beta$ with arbitrary $\beta\in H_C^u$.
Then $G^H\alpha$ is actual and
$A^HG^H\alpha=A^HG^HA^H\beta=A^H\beta=\alpha$.
No algebraicity of the originally chosen $\beta$ is used.
\end{proof}

\paragraph{Why this is a correspondence-class statement.}
The implication (PI20.6) to the correction formula uses the action
on full cohomology, including nonalgebraic classes.  An equality
of operators only on numerical groups is not enough.  Nor is
nilpotence of a numerically trivial error enough: in the abstract
algebra $\mathbf Q[\epsilon]/(\epsilon^2)$, modulo the ideal
$(\epsilon)$, the element $A=\epsilon$ has numerical image zero,
but $AGA=0\ne A$ for every $G$ in that algebra.  Its error is
nilpotent and nonzero.  This is only an algebraic limitation, not
a geometric example or a counterexample to G.

\subsection{A finite divisor certificate for a surface model}

\begin{theorem}[Surface incidence admits an algebraic cohomological inverse]
\label{r20pi:surface}
Suppose $d=2$.  Every choice of $G$ satisfying (PI20.4) satisfies
(PI20.6) for every prime $\ell\ne\operatorname{char}k$.
Consequently
\[
                         A^HG^HA^H=A^H
 \tag{PI20.13}\label{r20pi:surfaceinverse}
\]
on all geometric component cohomology.  The same actual rational
divisor matrix $G$ works for every such $\ell$.
\end{theorem}

\begin{proof}
Each product $D_j\times D_i$ is a smooth projective fourfold
which is a product of two arbitrary smooth projective surfaces.
The error $E_{ij}$ has codimension three, hence is an actual
rational one-cycle.  It is numerically zero by (PI20.4).
The frozen SP19 finite-field surface-product comparison proves
that its geometric $\ell$-adic class is zero, for every invertible
$\ell$.  Recall the precise classical reason behind that input:
Lefschetz $B$ is known for surfaces and stable under products;
its algebraic inverse in degree two sends a numerically zero
curve to a numerically zero divisor; divisor numerical and
homological equivalence coincide.  No Tate surjectivity for
surface $H^2$ or middle-degree Hodge positivity is assumed.
Thus every (PI20.6) holds, and the action identity follows.
\end{proof}

\begin{proposition}[A rational linear certificate]
\label{r20pi:certificate}
Still let $d=2$.  For every ordered pair $(k,l)$ choose actual
rational divisors
\[
 g_{kl,\mu}\in\operatorname{CH}^1(D_l\times D_k)_{\mathbf Q}
 \tag{PI20.14}\label{r20pi:basis}
\]
whose numerical classes form a supplied basis of
$N^1(D_l\times D_k)_{\mathbf Q}$.
For every ordered pair $(i,j)$ choose actual divisor representatives
$\eta_{ij,\nu}$ of such a basis on $D_j\times D_i$.
There is a solution over $\mathbf Q$ to the finite linear system
\[
 \sum_{k,l,\mu} M_{(i,j,\nu),(k,l,\mu)}\lambda_{kl,\mu}
                           =b_{i,j,\nu},
 \tag{PI20.15}\label{r20pi:linearsystem}
\]
where the entries are the native intersection numbers
\[
 \begin{split}
 M_{(i,j,\nu),(k,l,\mu)}
   &=\deg_{D_j\times D_i}
       \left((A_{ik}\circ g_{kl,\mu}\circ A_{lj})
                                      \cdot\eta_{ij,\nu}\right),\\
 b_{i,j,\nu}
   &=\deg_{D_j\times D_i}(A_{ij}\cdot\eta_{ij,\nu}).
 \end{split}
 \tag{PI20.16}\label{r20pi:matrix}
\]
Every rational solution, with
$G_{kl}=\sum_\mu\lambda_{kl,\mu}g_{kl,\mu}$, certifies
(PI20.13) for every invertible $\ell$.
The existence statement does not assert an algorithm that finds
complete Neron--Severi bases for arbitrary input surfaces.
\end{proposition}

\begin{proof}
The correspondence $A$ is fixed, so $AGA$ is linear in the
coefficients of $G$.  The entry
$(AGA-A)_{ij}$ is a curve class on the fourfold $D_j\times D_i$.
It is numerically zero exactly when it has zero intersection
with every numerical divisor class there.  The supplied basis
is therefore a complete set of tests, and expanding those
equalities gives (PI20.15)--(PI20.16).

Lemma~\ref{r20pi:numericalinverse} gives a numerical divisor
matrix $G$.  Expand it in the supplied rational bases; its
coefficients solve the linear system.  Conversely a solution
gives actual divisors and proves numerical triviality of every
error entry by the same perfect numerical pairing.  Apply
Theorem~\ref{r20pi:surface}.  The products and compositions in
(PI20.16) are actual Chow operations before degrees are taken.
\end{proof}

\subsection{The arbitrary-coefficient bridge and its scope}

\begin{corollary}[Tensoring a certified incidence correspondence]
\label{r20pi:coefficient}
Suppose (PI20.6) is verified at an invertible prime $\ell$.
Let $T/k$ be any smooth projective geometrically integral variety,
without a numerical--homological comparison assumption on $T$
or on $T\times D_i$.  On the component diagram
$D_i^T=T\times D_i$, with edges $T\times E_{ij}$, one has
\[
 A_T=\Delta_T\boxtimes A,\qquad
 G_T=\Delta_T\boxtimes G,\qquad
 A_T^HG_T^HA_T^H=A_T^H.
 \tag{PI20.17}\label{r20pi:tensor}
\]
Consequently all actual-span $A_T$-kernel defects and supported
CR19 defects vanish in every degree, including for mixed actual
cycle classes on the products.
\end{corollary}

\begin{proof}
Reorder the factors in the displayed exterior correspondences
to view them as correspondences on
$(T\times D_j)\times(T\times D_i)$.
The codimensions are respectively $\dim T+d+1$ and
$\dim T+d-1$, so their degrees are still $+1$ and $-1$.
Flat base change and the projection formula identify the actual
restriction and Gysin maps of this product diagram with
$\Delta_T\boxtimes A$.  Composition gives
\[
 A_TG_TA_T-A_T
              =\Delta_T\boxtimes(AGA-A)
\]
as actual correspondences.  Its geometric cohomology class is
zero by (PI20.6), compatibility with exterior products, and
Kunneth.  This proves the identity on all geometric cohomology.

The corrected-lift and actual-image arguments of
Corollary~\ref{r20pi:correctedlift} apply verbatim.  They apply
$G_T$ directly to the supplied actual cycle class and hence
do not require that class to be an exterior product.  Moreover
$A_TG_TA_T-A_T$ is numerically zero, either by the adequacy of
numerical equivalence for exterior products or by its zero
cohomology class, so the numerical part of the correction is
also valid.
\end{proof}

\paragraph{Native fields and final analytic use.}
All constructions above are native.  If finitely many components
are not geometrically integral, pass to one finite constant
extension that splits them, and retain the full degree in all
native norm and intersection formulas.  A native numerically
zero correspondence remains numerically zero after such an
extension: pair with any new complementary cycle and apply the
proper finite norm and projection formula.  Geometric cohomology
vanishing is then unchanged by that constant extension.

The arbitrary-coefficient bridge becomes a local Green theorem
only after its separate actual inertia-zero input and the
supplied model/alteration hypotheses have been checked.  The
AS20 application makes those checks for a good-reduction
coefficient and a surface factor, including all finite-extension
new cycles.  This section itself constructs no analytic pushforward
and no Green form by pushing along an arbitrary proper map.
Conversion from zero actual periods uses the frozen FD17/VT17
interfaces, preserving exact original support and all compact
ordinary actual tests.

\paragraph{What remains open in general dimension.}
Jannsen's theorem always gives (PI20.4), but in dimension $d>2$
its error has codimension $d+1$ on a $2d$-dimensional component
product.  The surface curve-comparison argument no longer
applies.  Vanishing of its geometric class is a stronger finite
sufficient input, not a necessary reformulation of G.  The
actual period image could vanish with a nonzero error action or
a nonzero $\Delta_A$.  No general vanishing, no universal
algebraic incidence inverse, and no new counterexample is
claimed here.

\paragraph{Dependencies and source record.}
The actual incidence signs, actual-span projections, and the
original-field Green bridge are frozen DQ17/ER17 and AK19/CR19.
The numerical generalized inverse uses U.~Jannsen,
\emph{Motives, numerical equivalence, and semi-simplicity},
Invent. Math.\ 107 (1992), 447--452, Theorem~1; the twist and
correspondence convention is stated on pp.~447--448 of that
primary paper.  Its numerical semisimplicity is not interpreted
as a cohomological comparison theorem.  The surface comparison
is SP19, ultimately using Kleiman's Proposition~4.3 for Lefschetz
$B$ of surfaces and their products.  All remaining assertions
are the actual Chow calculations and finite-dimensional arguments
written above.  No external priority claim is made.
